\section{Appendix}

\subsection{More Examples of Multi-Modal 3D Shape Retrieval}


In Figures~\ref{fig:sup_image_retrieval} and \ref{fig:sup_pc_retrieval}, we showcase more examples of multi-modal 3D shape retrieval.

\begin{figure}[h]
    \vspace{-0.5em}
    \centering
    \includegraphics[width=\linewidth]{figures/image_retrieval_supp.png}
    \caption{\textbf{Image-input 3D shape retrieval.} In each triplet, we present the input image and two 3D shapes retrieved using OpenShape embeddings from the Objaverse~\cite{objaverse} dataset. Input images are from \url{unsplash.com}.}
    \label{fig:sup_image_retrieval}
\end{figure}

\begin{figure}[h]
    \centering
    \includegraphics[width=\linewidth]{figures/pc_retrieval_supp.png}
    \caption{\textbf{Point cloud-input 3D shape retrieval.} In each triplet, we present the input point cloud and two 3D shapes retrieved using OpenShape embeddings from the Objaverse~\cite{objaverse} dataset.}
    \label{fig:sup_pc_retrieval}
\end{figure}

\subsection{More Examples of Shape-Conditioned Multimodal Generation}

In Figure~\ref{fig:sup_caption} and Figure~\ref{fig:sup_sd}, we showcase more examples of point cloud captioning and point cloud-conditioned image generation.

\begin{figure}[h]
    \centering
    \includegraphics[width=0.97\linewidth]{figures/supp_caption.png}
    \caption{\textbf{Point cloud captioning}. In each row, we show the input point clouds on the left and the generated captions on the right.}
    \label{fig:sup_caption}
\end{figure}

\begin{figure}[h]
    \centering
    \includegraphics[width=0.97\linewidth]{figures/supp_sd.jpg}
    \caption{\textbf{Point cloud-conditioned image generation}. Each row shows three examples (input point clouds and generated images).}
    \label{fig:sup_sd}
\end{figure}

\subsection{Details on Raw Text Generation and Filtering}

\subsubsection{Raw Text Generation}
We leverage the metadata from the four datasets to generate the raw texts. Although the original datasets may contain numerous attributes for each shape, we carefully choose the most informative ones to compose the text, ensuring its quality and relevance.

\noindent \textbf{Objaverse}:We utilize the \texttt{name} associated with each shape to serve as the text.

\noindent \textbf{ShapeNetCore}: For each shape, we generate three types of texts: (a) the \texttt{name}, (b) the \texttt{category name} (with a total of 55 categories), and (c) the concatenation of the \texttt{sub-category names} (with a total of 336 sub-categories), separated by commas.

\noindent \textbf{3DFuture}: For each shape, we generate two types of texts: (a) the \texttt{category}, and (b) the concatenation of \texttt{category}, \texttt{style}, \texttt{theme}, and \texttt{material},  separated by commas.

\noindent \textbf{ABO}:  
For each shape, we generate two types of texts: (a) the \texttt{item\_name}, and (b) the \texttt{product\_type}.

In this way, we generate one or more raw texts for each shape.

\subsubsection{Raw Text Filtering}

We employ GPT-4~\cite{openai2023gpt4} to filter out uninformative raw texts. To accomplish this, we divide all the raw texts into batches, each containing 256 entries, and process each batch independently using GPT-4. Here is an example illustrating the prompt we used and the corresponding response generated by GPT-4. 

\begin{minipage}[t]{11.5cm}
\begin{boxM}
I am analyzing a 3D dataset with various text descriptions for the 3D models. However, many of these texts are inaccurate or uninformative, and therefore, not suitable as descriptions for 3D models. I need your help to identify such incorrect texts. Specifically, if a text primarily consists of irrelevant or uninformative content, such as timestamps, model numbers, incomprehensible descriptions, random filenames (e.g., "my project"), random characters, etc., please respond with "N". If a text contains a clear noun (or noun phrase) that could potentially describe a 3D object, please respond with "Y". You will find a list of texts below, and each line contains a three-digit ID and associated text. For each text, please respond with "Y" or "N", following the ID number (e.g., "001 Y" or "002 N"). Please evaluate all 256 texts.

000 New project ( 19 )

001 3December - Chemestry

002 Fake Brand Soda Can

003 Spartan Shild

004 Apple3d

005 Landmine

006 FaunveinB-S

007 FIGURA 5

008 Sphero Blue

009 Sofa

010 Maddox

011 A3 Complete

012 Suspension Bridge

013 Maung

014 Captain-americas-shield

015 sphorb4

......
\end{boxM}
\end{minipage}
\hfill
\begin{minipage}[t]{2cm}
\begin{boxM}
000 N

001 Y

002 Y

003 Y

004 Y

005 Y

006 N

007 N

008 Y

009 Y

010 N

011 N

012 Y

013 N

014 Y

015 N

......

\end{boxM}
\end{minipage}

Afterwards, we combine all the responses to create the final filtering results, effectively removing approximately 30\% of the raw texts.

\subsection{Details on the Backbone Scaling Experiment}

In Figure 4 of the main paper, we investigate the performance and scalability of various backbones when scaling up their model sizes. For this experiment, we employ a default resolution of 10,000 points for input point clouds, a batch size of 200, and conduct the experiment on a single A100 GPU. In general, if instructions are given in the original paper of a backbone, we scale up the model as instructed. Otherwise, we scale up the model by expanding width or depth (i.e., stacking blocks or layers). Specifically, we scale up each backbone as follow:

\paragraph{PointBERT~\cite{yu2021pointbert}} The scaling parameters are shown in Table~\ref{tab:pbertscl}. We scaled PointBERT to 72.1M parameters beyond the 32.3M version reported in Figure 4 of the main paper. However, at this scale, the model dramatically overfits on the training data and performs worse on all benchmarks than the 32.3M version.

\begin{table}[h]
\centering
\caption{Hyperparameters for scaling up PointBERT~\cite{yu2021pointbert}.}
\label{tab:pbertscl}
\begin{tabular}{ccccccc}
\toprule
\# Parameters & \# Layers & Width & \# Heads & MLP Dim & \# Patches & Patch Embed Dim  \\ 
\midrule
5.1M          & 6         & 256   & 4        & 1024    & 64         & 96                   \\
13.3M         & 6         & 512   & 8        & 1024    & 64         & 128                  \\
32.3M         & 12        & 512   & 8        & 1536    & 384        & 256                  \\
\color{gray}72.1M         & \color{gray}12        & \color{gray}768   & \color{gray}12        & \color{gray}2304    & \color{gray}512        & \color{gray}256                  \\
\bottomrule
\end{tabular}
\end{table}

\paragraph{SparseConv~\cite{choy20194d}} The smallest version (5.3M parameters) of the model is adapted from the MinkowskiFCNN model by adjusting the width of the final convolution and linear layers. The remaining three models are adaptations of MinkowskiResNet, each varying in the number of basic ResNet blocks used. See Table~\ref{tab:sparse_cov_scaling} for the specific scaling parameters.

\begin{table}[h]
  \centering
  \caption{Hyperparameters for scaling up SparseConv~\cite{choy20194d}.}
    \begin{tabular}{ccc}
    \toprule
    \# Parameters & \# Convolution Layers & \# Linear Layers \\
    \midrule
    5.3M  & 7     & 4 \\
    29.0M & 18    & 3 \\
    33.7M & 26    & 3 \\
    41.3M & 42    & 3 \\
    \bottomrule
    \end{tabular}%
  \label{tab:sparse_cov_scaling}%
\end{table}%

\paragraph{PointNeXt~\cite{qian2022pointnext}} PointNeXt is proposed as a scalable version of PointNet++~\cite{qi2017pointnet2}, and includes S/B/L/XL variants in the original paper. We simply adopt these official configurations.

\paragraph{DGCNN~\cite{wang2019dynamic} and PointNet~\cite{qi2017pointnet}} For these two backbones without a hierarchical structure, we increase the width of each layer proportionally to scale up to 4xPointNet and 2xDGCNN before we hit the GPU memory limit. As the models operate completely on dense points, it is impractical to use the default 10k-point resolution. We thus reduce the input resolution for the two backbones, resulting in 1k points for DGCNN and 4k points for PointNet.

\subsection{Details on Training and Evaluation}

\paragraph{Training Details}

We freeze the CLIP text and image encoders and train the 3D encoder and two projection heads on our ensembled dataset using the cross-modal contrastive loss. We train the model on a single A100 GPU with a batch size of 200. Since we precache the text and image CLIP embeddings of all shapes, the training is greatly accelerated and takes about 300 A100 hours for convergence. We utilize an exponential learning rate schedule, and employ an range test to find the initial learning rate. For 32.3M version of PointBERT, we utilize a learning rate of $5e-4$; for 72.1M version of PointBERT, we utilize a learning rate of $4e-4$; and for other models, we utilize a learning rate of $1e-3$. For hard-negative mining, the number of seed shapes $s$ is set to 40, and the number of neighbors $m$ is set to 5 per shape, and the threshold $\delta$ is set to 0.1. 

\paragraph{Fine-tuning CLIP Text and Image Encoders?}

After training OpenShape-PointBERT, we conducted experiments to unfreeze and finetune the CLIP text encoder for a single epoch. However, the results obtained did not demonstrate any noticeable improvement on the benchmarks. Moreover, we observed that finetuning the CLIP text encoder could potentially undermine the generalization capabilities of CLIP and hinder the integration of OpenShape embeddings into existing CLIP-based models. As a result, we choose to freeze the CLIP encoders throughout the entire training process.

\paragraph{Evaluation Details} We evaluated all baselines using their publicly released pretrained checkpoints. Additionally, we retrained ULIP~\cite{xue2022ulip} on our ensembled training shapes using their official code base and backbone networks. Note that the retrained ULIP model utilized the original raw texts from the four datasets during training (prompt engineering is also applied), rather than our filtered and enriched texts. For ModelNet40~\cite{wu20153d}, the evaluation is conducted on the test split with 2,468 shapes. Regarding ScanObjectNN~\cite{uy2019revisiting}, we follow ULIP~\cite{xue2022ulip} to evaluate on the OBJ\_ONLY version, which contains 581 test shapes. For Objaverse-LVIS~\cite{objaverse}, the input is 10,000 sampled points with point colors. For ModelNet40~\cite{wu20153d}, the input is 10,000 sampled points without color. For ScanObjectNN~\cite{uy2019revisiting}, we utilize the official 2,048 points without color as input. All methods use the same input during evaluation. The forward inference time on an A100 GPU for a 10,000-point point cloud is approximately 0.9ms for OpenShape-SparseConv and 3.8ms for OpenShape-PointBERT.


\subsection{Details on Shape-Conditioned Multimodal Generation}
\paragraph{Point Cloud Captioning} 

CLIPCap~\cite{mokady2021clipcap} utilizes a 10-token prefix generated from CLIP image embeddings to enable GPT-2 for captioning. In order to align with the off-the-shelf CLIPCap model, we trained a variant of OpenShape-PointBERT that employs CLIP ViT-B/32 embeddings instead of OpenCLIP ViT-G/14 used in other experiments. Consequently, we directly input the point cloud encoding, \emph{without normalization}, into CLIPCap for captioning.

\paragraph{Point Cloud Conditioned Image Generation} We take the Stable Diffusion v2.1 unCLIP model~\cite{ramesh2022hierarchical} for image generation and replace the CLIP image condition encoder with our OpenShape encoder to perform image generation conditioned on point clouds (and optionally text prompts). The unCLIP model takes CLIP ViT-L/14 embeddings without normalization as input. To match the embedding space, we trained a variant of OpenShape-PointBERT with CLIP ViT-L/14 embeddings. Additionally, we noticed a significant mismatching of scales ($L_2$-norm of embedding vectors) between ViT-L/14 image embeddings and OpenShape embeddings. To mitigate this issue, we perform a re-normalization on OpenShape embeddings to a $L_2$-norm of $\frac{1}{2}\sqrt{768}$, which is our observed mean $L_2$-norm of ViT-L/14 image embeddings. We use 50 diffusion steps. The guidance scale can be tuned freely.



